\section{群作用}
\label{sec:lie-actions}

上一节的群只规定了变换之间怎样复合。同一个旋转可以转动一个点、一整组点，也可以转动一幅定义在平面上的图像；被变换的对象不同，计算公式也不同。本节先规定群元素怎样作用在给定集合上，并核对作用与复合相容；再用轨道与稳定子描述对象之间通过变换建立的关系；最后在输入与输出的作用都给定后，说明映射的结果应怎样随输入变换。

\begin{definition}[群作用]
\label{def:lie-group-action}
设$G$是群，$X$是集合。映射$G\times X\to X$，$(g,x)\mapsto g\cdot x$若满足
\begin{equation}
 e\cdot x=x,\qquad (g_1g_2)\cdot x=g_1\cdot(g_2\cdot x)\qquad(g_1,g_2\in G,\ x\in X),
 \label{eq:lie-group-action}
\end{equation}
就称为$G$在$X$上的\term{群作用}（group action）。
\end{definition}

第一条说单位元什么也不做，第二条说先用$g_2$作用、再用$g_1$作用，与用乘积$g_1g_2$作用一次相同。对固定的$g$，映射$x\mapsto g\cdot x$有逆映射$x\mapsto g^{-1}\cdot x$，因为$g^{-1}\cdot(g\cdot x)=(g^{-1}g)\cdot x=x$。所以群作用相当于一个同态，把每个$g$对应为$X$上的一个双射。

\paragraph{同一个群在不同对象上的作用}
$\mathrm{SO}(2)$在$\mathbb R^2$上按矩阵乘法$\bm R\cdot\bm x=\bm R\bm x$作用，两条规则就是$\bm I\bm x=\bm x$与矩阵乘法的结合律。$\mathbb R^2$也按$\bm a\cdot\bm x=\bm x+\bm a$作用在自身上。

点云由$n$个点组成，记为$\bm X\in\mathbb R^{n\times D}$，第$i$行$\bm X_{i,:}$是第$i$个点的坐标。$S_n$按$\pi\cdot\bm X=\bm P_\pi\bm X$作用，其中$\bm P_\pi$沿用上一节的置换矩阵；由$\bm P_\pi\bm e_j=\bm e_{\pi(j)}$可以算出$(\bm P_\pi\bm X)_{i,:}=\bm X_{\pi^{-1}(i),:}$，即原来的第$j$行被搬到第$\pi(j)$行。$\mathrm{SO}(D)$则按$\bm R\cdot\bm X=\bm X\bm R^\top$作用，把每一行$\bm x^\top$变成$(\bm R\bm x)^\top$。两种作用都写成矩阵乘法，前者改变点的编号，后者移动点的位置，被改变的对象和索引完全不同。

旋转还可以作用在平面图像上。把图像记为函数$f:\mathbb R^2\to\mathbb R$，令
\[
(\bm R\cdot f)(\bm x)=f(\bm R^{-1}\bm x).
\]
这里必须用$\bm R^{-1}$：新图像在$\bm x$处的值，来自旧图像中转到$\bm x$之前的位置。验证第二条规则：$(\bm R_1\cdot(\bm R_2\cdot f))(\bm x)=(\bm R_2\cdot f)(\bm R_1^{-1}\bm x)=f(\bm R_2^{-1}\bm R_1^{-1}\bm x)=f((\bm R_1\bm R_2)^{-1}\bm x)$，恰为$((\bm R_1\bm R_2)\cdot f)(\bm x)$。若改写成$f(\bm R\bm x)$，复合次序就会颠倒，在非交换群上不再满足式~\eqref{eq:lie-group-action}。

\subsection{轨道与稳定子}
\label{sec:lie-orbits-stabilizers}

给定作用后，可以问两个方向相反的问题：从一个对象出发，经过群中的变换能到达哪些对象；哪些变换根本不改变这个对象。前者给出轨道，后者给出稳定子。

\begin{definition}[轨道与稳定子]
\label{def:lie-orbit-stabilizer}
设群$G$作用于$X$，$x\in X$。称
\[
G\cdot x=\{g\cdot x:g\in G\},\qquad G_x=\{g\in G:g\cdot x=x\}
\]
分别为$x$的\term{轨道}（orbit）与\term{稳定子}（stabilizer）。若$X$只有一条轨道，就称作用是\term{传递作用}（transitive action）。
\end{definition}

轨道是$X$的子集，稳定子是$G$的子集，二者不在同一个空间里。规定$x\sim y$当且仅当存在$g\in G$使$y=g\cdot x$。单位元给出$x\sim x$；若$y=g\cdot x$，则$x=g^{-1}\cdot y$，给出对称性；若$y=g\cdot x$、$z=h\cdot y$，则$z=(hg)\cdot x$，给出传递性。所以$\sim$是等价关系，其等价类恰是轨道。由等价关系与划分的对应（可参考\BookChapterRef{chap:sets-relations-maps}{卷一第1章“集合、关系与映射”}），不同的轨道互不相交，且所有轨道的并是$X$。轨道组成的商集记为$X/G$。

稳定子则总是子群。若$g,h\in G_x$，则$h^{-1}\cdot x=h^{-1}\cdot(h\cdot x)=x$，于是$(gh^{-1})\cdot x=g\cdot(h^{-1}\cdot x)=x$，由引理~\ref{lem:lie-subgroup-criterion}得$G_x$是子群。

\paragraph{旋转作用的轨道}
$\mathrm{SO}(2)$作用于$\mathbb R^2$时，$\bm x\neq\bm0$的轨道是半径为$\lVert\bm x\rVert_2$的圆：旋转保持长度，并且同一圆上任意两点之间都有一个旋转把一点送到另一点。原点的轨道只有原点自身。非零点的稳定子只含$\bm I$，原点的稳定子是整个$\mathrm{SO}(2)$。轨道可以用半径$r\geq0$编号，所以$\mathbb R^2/\mathrm{SO}(2)$可以识别为$[0,\infty)$，原点的轨道对应端点$0$。

$\mathrm{SO}(3)$作用于$\mathbb R^3$时，非零点的轨道是球面。以$\bm e_3=(0,0,1)^\top$为例，它的轨道是单位球面$S^2$，稳定子是绕第三坐标轴的全部旋转
\[
\left\{\begin{bmatrix}\bm R(\theta)&\bm0\\\bm0^\top&1\end{bmatrix}:\theta\in\mathbb R\right\},
\]
这是一个与$\mathrm{SO}(2)$同构的子群。这里有三个不同的空间：群$\mathrm{SO}(3)$由旋转组成，单个对象$\bm e_3$是一个点，轨道$S^2$是一个曲面。许多不同的旋转把$\bm e_3$送到同一个点，所以不能把$\mathrm{SO}(3)$与球面$S^2$等同；第~\ref{sec:lie-quotients}节会说明$S^2$正是$\mathrm{SO}(3)$按这个稳定子取商的结果。

\paragraph{轨道与陪集}
两个变换把$x$送到同一个点，当且仅当它们只差一个稳定子中的元素：$g\cdot x=h\cdot x$等价于$(h^{-1}g)\cdot x=x$，即$h^{-1}g\in G_x$。为了记录“只差$G_x$中元素”的变换，引入下面的集合。

\begin{definition}[陪集]
\label{def:lie-coset}
设$H$是群$G$的子群，$g\in G$。称$gH=\{gh:h\in H\}$为$H$的一个左\term{陪集}（coset）。全体左陪集组成的集合记为$G/H$。
\end{definition}

$h^{-1}g\in G_x$恰等价于$g\in hG_x$，于是$gG_x=hG_x$当且仅当$g\cdot x=h\cdot x$。这说明映射$gG_x\mapsto g\cdot x$不依赖代表元，并且是从$G/G_x$到轨道$G\cdot x$的双射。

\begin{theorem}[轨道—稳定子关系]
\label{thm:lie-orbit-stabilizer}
设群$G$作用于$X$，$x\in X$。则$gG_x\mapsto g\cdot x$是$G/G_x$到$G\cdot x$的双射。若$G$是有限群，则$\lvert G\cdot x\rvert\,\lvert G_x\rvert=\lvert G\rvert$。
\end{theorem}

\begin{proof}
双射已在上一段证明。有限情形下，左乘$g$是$G$到自身的双射，把$G_x$一一地映成$gG_x$，所以每个陪集都有$\lvert G_x\rvert$个元素。不同陪集互不相交：若$k\in gG_x\cap hG_x$，则$k\cdot x=g\cdot x=h\cdot x$，从而$gG_x=hG_x$。每个$g$都在自己的陪集$gG_x$中，所以$G$被分成$\lvert G/G_x\rvert=\lvert G\cdot x\rvert$个大小为$\lvert G_x\rvert$的块。
\end{proof}

用$S_3$作用于$\{1,2,3\}$核对：$1$的轨道是$\{1,2,3\}$，稳定子是$\{\operatorname{id},\tau\}$，其中$\tau$交换$2$与$3$，确有$3\times2=6$。对点云，若$\bm X$的$n$行互不相同，则只有恒等置换保持$\bm X$，轨道有$n!$个元素，对应同一组点的$n!$种编号方式；若恰有两行相同，交换这两行也保持$\bm X$，稳定子有$2$个元素，轨道只有$n!/2$个不同的矩阵。

\subsection{不变映射与等变映射}
\label{sec:lie-invariant-equivariant}

在点云上计算一个量时，任务往往对变换有明确要求。点的重新编号不应改变“这组点的重心在哪里”；整体旋转之后，重心应跟着转到新的位置，而点之间的距离则不应改变。两种要求都是在比较“变换后再计算”与“计算后再变换”，区别在于输出是否也受到作用。

\begin{definition}[不变映射与等变映射]
\label{def:lie-invariant-equivariant}
设群$G$作用于$X$与$Y$，映射$f:X\to Z$，$F:X\to Y$。若
\begin{equation}
 f(g\cdot x)=f(x)\qquad(g\in G,\ x\in X),
 \label{eq:lie-invariant}
\end{equation}
就称$f$是\term{不变映射}（invariant map）。若
\begin{equation}
 F(g\cdot x)=g\cdot F(x)\qquad(g\in G,\ x\in X),
 \label{eq:lie-equivariant}
\end{equation}
就称$F$是\term{等变映射}（equivariant map），其中左边的点是$X$上的作用，右边的点是$Y$上的作用。
\end{definition}

不变性只需要输入的作用；等变性必须同时给定输入与输出的作用，否则式~\eqref{eq:lie-equivariant}没有意义。若$Y$上的作用是平凡作用$g\cdot y=y$，等变就退化为不变。所以不变是等变的特例，而不是另一种更强或更好的要求。

\paragraph{同一次变换下的核对}
取$\bm x=(1.3,0.5)^\top$，绕原点逆时针转$90^\circ$，$\bm R(\pi/2)\bm x=(-0.5,1.3)^\top$。长度$f(\bm x)=\lVert\bm x\rVert_2$在变换前后都等于$\sqrt{1.94}$，满足式~\eqref{eq:lie-invariant}。恒等映射$F(\bm x)=\bm x$满足$F(\bm R\bm x)=\bm R\bm x=\bm R F(\bm x)$，是等变映射，但$F(\bm R\bm x)=(-0.5,1.3)^\top\neq F(\bm x)$，所以它不是不变映射。等变保留的是“输出随输入同步变化”的规则，而不是输出保持不动。

\paragraph{点云上的例子}
对点云$\bm X\in\mathbb R^{n\times D}$，令$\bm 1\in\mathbb R^n$为全一向量，重心为$\bm m(\bm X)=\frac1n\bm X^\top\bm 1\in\mathbb R^D$。重新编号时，$\bm m(\bm P_\pi\bm X)=\frac1n\bm X^\top\bm P_\pi^\top\bm 1=\bm m(\bm X)$，因为置换矩阵的每一行恰有一个$1$，$\bm P_\pi^\top\bm 1=\bm 1$。旋转时，$\bm m(\bm X\bm R^\top)=\frac1n\bm R\bm X^\top\bm 1=\bm R\,\bm m(\bm X)$。所以重心对置换不变，对旋转等变，输出的旋转作用是$\bm y\mapsto\bm R\bm y$。

两两距离矩阵$\bm D(\bm X)_{i,j}=\lVert\bm X_{i,:}-\bm X_{j,:}\rVert_2$对旋转不变，因为旋转保持差向量的长度。它对置换却不是不变的：重新编号后，$\bm D(\bm P_\pi\bm X)$的$(i,j)$元是原来第$\pi^{-1}(i)$、$\pi^{-1}(j)$两点的距离，即$\bm D(\bm P_\pi\bm X)=\bm P_\pi\bm D(\bm X)\bm P_\pi^\top$。若在矩阵空间上规定作用$\pi\cdot\bm M=\bm P_\pi\bm M\bm P_\pi^\top$，$\bm D$就是等变的。同一个映射是否“对置换不变”，取决于问的是哪一种输出作用。

\paragraph{不变映射与轨道}
不变映射在同一条轨道上取相同的值，所以它只依赖输入所在的轨道。记$q:X\to X/G$为$q(x)=G\cdot x$，令$\widetilde f(G\cdot x)=f(x)$；若$G\cdot x=G\cdot y$，则$y=g\cdot x$，由不变性$f(y)=f(x)$，所以$\widetilde f$不依赖代表元，并且$f=\widetilde f\circ q$。不变映射正是忽略了同一轨道内的差别。

这种忽略可能过多。$f$可以把不同的轨道映到同一个值，例如把每个点云都映到零的常值映射对一切作用都不变，却不区分任何点云。不变性只说明$f$没有利用轨道内的差别，不说明$f$保留了区分不同轨道所需的信息。

\paragraph{复合与汇总}
等变映射之后接不变映射，结果不变：若$F$等变、$f$在$Y$上不变，则$f(F(g\cdot x))=f(g\cdot F(x))=f(F(x))$。两个等变映射的复合仍等变，理由相同。例如对每个点使用同一个映射$\phi:\mathbb R^D\to\mathbb R^r$，令$\Phi(\bm X)_{i,:}=\phi(\bm X_{i,:})$，则$\Phi(\bm P_\pi\bm X)_{i,:}=\phi(\bm X_{\pi^{-1}(i),:})=(\bm P_\pi\Phi(\bm X))_{i,:}$，所以$\Phi$对置换等变；再对行求和，重新编号只改变求和次序，得到不变的$\sum_i\phi(\bm X_{i,:})$。反过来，若中间已经求和，丢掉了每个点的编号，后面就不能再恢复逐点的等变输出。
